\chapter{Microdochectomy}

\section*{Introduction \& Clinical Indications}
Microdochectomy is a targeted surgical procedure that involves the excision of a single breast duct responsible for pathologic nipple discharge. This minimally invasive intervention is typically performed within the subareolar breast tissue through a small, cosmetically favorable incision at the edge of the areola. The primary clinical scenario for its application is persistent, spontaneous, and unilateral nipple discharge originating from a single duct, particularly when the discharge is bloody, serosanguinous, or clear and watery. The procedure serves a dual purpose: to definitively diagnose the underlying cause of the discharge by providing tissue for histopathological examination and to therapeutically resolve the bothersome symptom. Microdochectomy is crucial for distinguishing between benign intraductal lesions, such as papillomas, and more concerning pathologies like atypical hyperplasia or rare forms of ductal carcinoma, thereby guiding appropriate subsequent management.

\begin{itemize}
  \item Single-Duct Excision
  \item Subareolar Duct Excision
  \item Terminal Duct Excision
  \item Nipple Duct Excision
  \item Microductectomy
  \item Major Duct Excision (when referring to the main lactiferous duct)
\end{itemize}

The fundamental principle of microdochectomy lies in the precise identification and complete removal of the specific breast duct that is the source of pathologic nipple discharge. This discharge is most commonly caused by an intraductal lesion, such as a benign papilloma, which is a small, wart-like growth within the ductal lumen. Other potential causes include duct ectasia, usual ductal hyperplasia, or, less frequently, atypical hyperplasia or carcinoma. 

By excising only the affected duct and its contents, the procedure aims to achieve both diagnostic clarity and therapeutic resolution while preserving as much healthy breast tissue as possible and minimizing cosmetic impact. The surgical rationale is based on the understanding that single-duct pathologic discharge is typically localized to a specific segment of the ductal system. Removing this segment ensures that the entire lesion, which may be too small to be reliably detected by imaging or cytology alone, is obtained for definitive histopathological analysis. 

This microscopic examination is critical for determining the nature of the lesion (benign, atypical, or malignant) and guiding any subsequent management. The technical goals include meticulous identification of the culprit duct, often aided by pre-operative localization techniques such as ductography or methylene blue injection; complete excision of the ductal segment containing the lesion; achieving adequate hemostasis to prevent hematoma formation; and ensuring a favorable cosmetic outcome by utilizing a circumareolar incision.

The understanding of nipple discharge and its underlying pathologies has evolved significantly over centuries. Early descriptions of breast conditions, including discharge, can be traced back to ancient Egyptian and Greek texts. The recognition of intraductal papillomas as a distinct cause of bloody nipple discharge dates back to the mid-nineteenth century, with detailed pathological descriptions by surgeons like Sir Astley Cooper (1829) and later by Billroth (1880s). 

Initially, surgical approaches to nipple discharge were often extensive and non-specific, involving large resections of breast tissue, such as Hadfield's procedure (total duct excision) or wide cone excisions, which frequently resulted in significant cosmetic deformity and loss of breast volume. The refinement of surgical techniques in the early to mid-twentieth century led to a paradigm shift towards more targeted excisions. 

The development of methods for precise duct identification, such as the use of fine probes or the injection of vital dyes like methylene blue into the discharging duct, allowed surgeons to isolate and remove only the affected duct. This evolution culminated in the modern microdochectomy, which prioritizes accurate diagnosis and symptom resolution through minimal intervention, thereby preserving breast cosmesis and function. Contemporary practice continues to emphasize precision, guided by advanced imaging (e.g., ductography, MRI) and a multidisciplinary approach to breast care, further enhancing diagnostic accuracy and patient outcomes.

Microdochectomy is indicated for specific clinical scenarios related to pathologic nipple discharge, particularly when a definitive diagnosis and symptom resolution are required.

\begin{itemize}
  \item Pathologic Nipple Discharge: When discharge is spontaneous, persistent, unilateral, and originates from a single duct, especially if it is bothersome to the patient.
  \item Bloody or Serosanguinous Discharge: This type of discharge carries a higher suspicion for underlying pathology, including malignancy, making excision imperative for definitive diagnosis.
  \item Clear or Watery Persistent Discharge: While less commonly associated with malignancy than bloody discharge, persistent clear or watery discharge from a single duct warrants investigation and often excision, particularly if other workup is inconclusive.
  \item Intraductal Lesion on Imaging: If imaging studies (e.g., ductography, breast MRI, ultrasound) identify an intraductal mass or abnormality within a single duct that correlates with the discharge.
  \item Persistent Discharge with Negative or Inconclusive Imaging: When initial imaging (mammography, ultrasound) is negative or inconclusive, but the pathologic discharge continues, microdochectomy provides a definitive tissue diagnosis.
  \item Atypical or Suspicious Cytology: If cytology performed on the nipple discharge fluid reveals atypical cells or other suspicious findings, surgical excision is necessary for accurate diagnosis and to rule out malignancy.
  \item Patient Anxiety and Preference: For patients with persistent, bothersome discharge who desire definitive diagnosis and treatment, even if initial workup is equivocal.
  \item Exclusion of Malignancy: In patients with risk factors for breast cancer, microdochectomy can help exclude malignancy as the cause of the discharge, providing reassurance or guiding further treatment.
\end{itemize}

Microdochectomy serves a dual role in the management of nipple discharge, acting as both a primary surgical treatment and a secondary diagnostic step. It is considered a primary surgical treatment because, once indicated, it directly addresses and resolves the symptom by removing the source of the discharge. For instance, if a benign intraductal papilloma is causing the discharge, its excision is often curative. The procedure is definitive in its therapeutic intent for the symptom of discharge.

However, in the broader diagnostic pathway for nipple discharge, microdochectomy is typically a secondary step. Initial evaluation always involves a thorough history and physical examination, followed by non-invasive imaging studies such as bilateral mammography and targeted breast ultrasound. Depending on the clinical picture, further imaging like breast MRI or ductography may be performed to localize the discharging duct or identify intraductal lesions. Nipple discharge cytology may also be obtained, though its sensitivity for malignancy is limited. 

Microdochectomy is usually reserved for cases where these initial investigations confirm single-duct pathologic discharge, especially if it is bloody, persistent, or associated with suspicious findings, and when a definitive tissue diagnosis is required that cannot be obtained through less invasive means. Thus, it is a primary intervention for a specific problem, but often follows a sequence of initial diagnostic evaluations to confirm its necessity and target.

\section*{Relevant Surgical Anatomy}
The breast is a modified apocrine gland composed of 15-20 lobes, each drained by a lactiferous duct that converges towards the nipple. These ducts widen beneath the areola to form lactiferous sinuses before narrowing again to open onto the nipple surface. Microdochectomy specifically targets one of these major lactiferous ducts and its associated branching system within the subareolar region. The nipple-areolar complex is richly innervated by branches of the 4th, 5th, and 6th intercostal nerves, which supply sensation to the skin and areola; careful dissection is paramount to preserve these nerves. 

The arterial supply to the breast is primarily from the internal mammary artery (medial aspect), lateral thoracic artery (lateral aspect), and intercostal arteries. Venous drainage largely parallels the arterial supply, emptying into the axillary and internal mammary veins. Lymphatic drainage is extensive, predominantly to the axillary lymph nodes (75\%), but also to the internal mammary nodes (20\%) and, less commonly, supraclavicular nodes. 

Surgical landmarks include the nipple itself, the areola, and the underlying subareolar fat pad where the major ducts reside. The procedure focuses on the terminal portion of the ductal system, close to the nipple, where most intraductal lesions causing discharge are located. 

\begin{itemize}
  \item Lactiferous ducts: Major ducts that transport milk from the lobules to the nipple.
  \item Nipple-areolar complex: The area encompassing the nipple and surrounding areola, rich in sensory nerves.
  \item Arterial supply: Internal mammary artery, lateral thoracic artery, and intercostal arteries.
  \item Venous drainage: Axillary and internal mammary veins.
  \item Lymphatic drainage: Primarily to axillary lymph nodes.
\end{itemize}

\section*{Preoperative Workup \& Preparation}
Thorough pre-operative preparation is essential to ensure patient safety and optimize surgical outcomes for microdochectomy. This typically begins with a comprehensive medical evaluation to assess the patient's overall health and fitness for surgery.

\begin{itemize}
  \item Comprehensive Medical Evaluation: This includes a detailed history and physical examination, review of current medications, assessment of any allergies, and evaluation of cardiac and pulmonary status.
  \item Laboratory Tests: Routine blood work typically includes a complete blood count (CBC), basic metabolic panel, and coagulation studies (PT/INR, PTT) to assess overall health, kidney function, and bleeding risk.
  \item Imaging Studies: Bilateral mammography and targeted breast ultrasound are standard to evaluate the breast tissue and help localize the discharging duct or any associated lesions. In some cases, ductography (galactography) or breast MRI may be performed for more precise localization or to rule out multifocal disease.
  \item Anesthesia Consultation: A pre-anesthetic evaluation is conducted to assess the patient's fitness for general anesthesia and discuss the anesthetic plan, including potential risks and alternatives.
  \item Medication Adjustments: Patients on blood-thinning medications (e.g., aspirin, NSAIDs, warfarin, novel oral anticoagulants) will receive specific instructions from their surgeon regarding when to discontinue these medications prior to surgery to minimize bleeding risk.
  \item NPO Status: Patients are instructed to fast (Nil Per Os) for a specified period before surgery, typically 6-8 hours for solids and 2 hours for clear liquids, to prevent aspiration during anesthesia.
  \item Antibiotic Prophylaxis: Intravenous prophylactic antibiotics, usually a first- or second-generation cephalosporin, may be administered shortly before the incision to reduce the risk of surgical site infection.
  \item Informed Consent: The surgeon will thoroughly discuss the procedure, its potential benefits, risks, alternatives, and expected recovery with the patient, obtaining informed consent for the surgery and anesthesia.
  \item Logistical Arrangements: Patients should arrange for a responsible adult to drive them home after the procedure and to provide support during the initial recovery period, as they will not be able to drive themselves due to the effects of anesthesia.
\end{itemize}

While microdochectomy is a safe and effective procedure, certain conditions may contraindicate or necessitate precautions before proceeding.

\textbf{Absolute Contraindications:}
\begin{itemize}
  \item Active Breast Infection: Presence of acute mastitis or a breast abscess in the operative field, as this significantly increases the risk of surgical site infection, complicates healing, and can obscure anatomical planes.
  \item Uncorrected Coagulopathy: Significant bleeding disorders that cannot be adequately managed pre-operatively, posing an unacceptable risk of hemorrhage during or after the procedure.
  \item Discharge from Multiple Ducts: Microdochectomy is designed for single-duct discharge. If discharge originates from multiple ducts, a total duct excision (Hadfield's procedure) or other diagnostic approaches may be more appropriate to address the diffuse nature of the pathology.
  \item Patient Refusal: The patient's informed decision not to proceed with surgery after a thorough discussion of risks, benefits, and alternatives.
\end{itemize}

\textbf{Relative Contraindications and Precautions:}
\begin{itemize}
  \item Pregnancy and Lactation: Nipple discharge is often physiologic during these periods. Surgery is generally deferred unless there is a high suspicion of malignancy, as it can interfere with lactation and carries potential risks to the fetus or infant.
  \item Severe Comorbidities: Patients with significant cardiac, pulmonary, or other systemic diseases that make general anesthesia or surgery itself high-risk may require extensive medical optimization or consideration of alternative diagnostic strategies if the risk outweighs the benefit.
  \item Extensive Prior Breast Surgery: Previous surgeries in the subareolar area can alter anatomy, making duct identification and dissection more challenging, potentially increasing operative time and risk of complications.
  \item Unclear Duct Localization: If, despite all pre-operative attempts (e.g., dye injection, probing, advanced imaging), the discharging duct cannot be reliably identified, the success of a targeted microdochectomy is compromised. In such cases, a more extensive central duct excision might be considered, or the procedure may be postponed until localization is possible.
  \item Imaging Suggestive of Extensive Malignancy: If pre-operative imaging strongly suggests an extensive malignancy (e.g., large mass, multifocal disease), a microdochectomy might be insufficient for complete removal, and a wider excision, lumpectomy, or mastectomy may be indicated as the primary treatment.
\end{itemize}

\section*{Surgical Technique \& Procedure}
Microdochectomy is typically performed as an outpatient procedure under general anesthesia, though local anesthesia with intravenous sedation may be used in select cases. The patient is positioned supine with the ipsilateral arm abducted to allow optimal access to the breast. The surgical field is prepped and draped in a sterile fashion.

A small, curvilinear incision, usually 2-3 cm in length, is made at the edge of the areola (circumareolar or periareolar incision) to ensure an optimal cosmetic outcome and minimize visible scarring. The key to the procedure is the accurate identification of the discharging duct. This is often facilitated by pre-operative injection of a small amount of methylene blue dye into the discharging duct, which stains the affected duct blue, or by inserting a fine lacrimal probe into the duct opening at the nipple. Gentle pressure on the breast may also help express discharge, guiding the surgeon. 

Subareolar dissection is then carefully performed to isolate the identified duct, often using fine dissecting scissors or electrocautery. Once isolated, the duct is ligated (tied off) proximally, close to the nipple, and distally, deeper within the breast tissue, to ensure complete removal of the affected segment and any contained lesion. The segment of the duct, along with its contents, is then excised. Meticulous hemostasis is achieved using electrocautery to prevent hematoma formation. The excised specimen is oriented (e.g., with sutures or ink) and sent for histopathological examination. The subcutaneous tissue is closed with fine absorbable sutures, and the skin is closed with absorbable or non-absorbable sutures, often using a subcuticular technique to further enhance cosmesis. A sterile dressing is applied, and drains are rarely necessary.

\section*{Complications \& Risks}
As with any surgical procedure, microdochectomy carries potential risks, which can be broadly categorized into general surgical complications and those specific to the breast and nipple-areolar complex.

\textbf{General Surgical Risks:}
\begin{itemize}
  \item Bleeding and Hematoma: Accumulation of blood under the skin, potentially causing swelling, pain, and bruising. A significant hematoma may require drainage and can increase the risk of infection.
  \item Infection: Surgical site infection (cellulitis or abscess) can occur, requiring antibiotics or, rarely, surgical drainage.
  \item Anesthesia Complications: Adverse reactions to anesthetic agents, including nausea, vomiting, allergic reactions, or more serious respiratory or cardiac complications, though rare.
  \item Pain: Post-operative pain, usually mild to moderate, managed effectively with oral analgesics.
  \item Scarring: Formation of a visible scar at the incision site. While typically well-hidden at the areolar edge, hypertrophic or keloid scars can occur in susceptible individuals.
\end{itemize}

\textbf{Procedure-Specific Risks:}
\begin{itemize}
  \item Nipple or Areolar Numbness/Altered Sensation: Damage to superficial nerves during dissection can lead to temporary or, less commonly, permanent changes in sensation of the nipple or areola.
  \item Changes in Nipple Shape or Position: Slight asymmetry, retraction, or distortion of the nipple or areola can occur, though efforts are made to minimize this through careful surgical technique.
  \item Persistent Nipple Discharge: If the entire affected duct or the underlying lesion is not completely removed, the discharge may recur, necessitating further investigation or re-excision.
  \item Recurrence of Benign Lesions: While the excised lesion is removed, new papillomas or other benign lesions can develop in other ducts or even in the same duct if excision was incomplete.
  \item Milk Fistula: If performed during lactation, there is a rare risk of developing a persistent tract that leaks milk, which can be challenging to manage.
  \item Damage to Adjacent Ducts: Unintended injury to nearby healthy ducts, potentially affecting future lactation capacity, especially if multiple ducts are inadvertently involved.
  \item Cosmetic Dissatisfaction: Despite best efforts to achieve an aesthetic outcome, some patients may be unhappy with the appearance of the nipple-areolar complex post-surgery.
  \item False Negative Diagnosis: Although rare with proper technique and pathology review, there is a theoretical risk that a small or multifocal malignancy could be missed if the excised tissue does not fully represent the pathology, especially if the discharging duct was not accurately identified.
\end{itemize}

\section*{Postoperative Care \& Recovery}
Following microdochectomy, patients are typically transferred to the Post-Anesthesia Care Unit (PACU) for close monitoring as they recover from anesthesia. Pain is usually mild and can be effectively managed with over-the-counter pain relievers or prescribed oral analgesics. Most microdochectomies are performed as outpatient procedures, allowing patients to be discharged home within a few hours once they are stable, awake, and comfortable. A sterile dressing will be applied to the incision site, and patients will receive detailed instructions on wound care, activity restrictions, and warning signs.

During the first 24-48 hours, it is common to experience some mild pain, swelling, and bruising around the incision site and nipple-areolar complex. Patients are advised to avoid strenuous activities, heavy lifting, and direct trauma to the breast. The incision site should be kept clean and dry, and showering instructions will be provided by the surgical team, usually allowing gentle washing after 24-48 hours. 

Within the first 1-2 weeks, any non-absorbable sutures will typically be removed at a follow-up appointment. Bruising and swelling will gradually subside, and most patients can progressively return to their normal daily activities, including light exercise. The pathology results, which are crucial for determining the definitive diagnosis and guiding any further management, are usually available within one to two weeks, and these will be discussed in detail with the surgeon. Long-term recovery involves scar maturation, which can take several months, and a gradual return of nipple sensation, though some degree of altered sensation may persist. Full recovery and return to all activities are typically expected within 4-6 weeks.

During microdochectomy, the surgeon's intraoperative findings typically include the identification of the discharging duct, often highlighted by methylene blue dye or a lacrimal probe. The duct may appear dilated or contain visible intraluminal material. The surrounding breast tissue is usually normal, though some fibrocystic changes might be noted. The excised ductal segment is then sent for histopathological examination, which provides the definitive diagnosis.

The histopathological report is the most critical aspect of understanding the microdochectomy results. The vast majority of microdochectomy specimens reveal benign findings, with an intraductal papilloma being the most common diagnosis, accounting for 50-80\% of cases. These are benign epithelial proliferations within the duct. Other benign findings may include duct ectasia (widening of the duct with associated inflammation), usual ductal hyperplasia (an increase in the number of cells lining the duct without atypia), or fibrocystic changes. Less commonly, the pathology report may reveal more significant findings such as atypical ductal hyperplasia (ADH), atypical lobular hyperplasia (ALH), ductal carcinoma in situ (DCIS), or, rarely, invasive carcinoma. 

These findings indicate a higher risk for future breast cancer or the presence of early-stage cancer, necessitating further management. The pathologist will also assess the margins of excision to determine if the lesion has been completely removed.

A normal and successful postoperative course following microdochectomy is characterized by a smooth recovery with resolution of the primary symptom and minimal complications. Patients can expect mild to moderate pain at the incision site for the first few days, which is well-controlled with oral analgesics. Swelling and bruising around the nipple-areolar complex are common but typically resolve within one to two weeks. 

The primary and most anticipated outcome is the complete cessation of nipple discharge, which usually occurs immediately after the procedure. Wound healing progresses without infection or dehiscence, leading to a small, well-healed scar at the areolar edge. Most patients can gradually resume light daily activities within a few days and return to full normal activities, including exercise, within 2-4 weeks. Nipple sensation may be temporarily altered, but often returns to near-normal over several weeks to months. The definitive pathology report, confirming a benign diagnosis, provides significant reassurance and peace of mind, marking a successful resolution of the clinical problem.

Postoperative care and follow-up are crucial after microdochectomy to monitor healing, discuss pathology results, and ensure long-term breast health.

\begin{itemize}
  \item Post-operative Clinic Visit: Typically scheduled 1-2 weeks after surgery for a wound check, removal of any non-absorbable sutures, and a detailed discussion of the histopathology report.
  \item Pathology Discussion: The surgeon will explain the findings of the excised tissue, discuss the implications for future care, and address any concerns, especially if atypical or malignant cells were found.
  \item Routine Breast Cancer Screening: Patients should continue to adhere to age-appropriate and risk-factor-based breast cancer screening guidelines, which usually include annual mammography. Additional imaging like ultrasound or MRI may be recommended based on personal risk factors or specific pathology findings (e.g., ADH, DCIS).
  \item Self-Breast Awareness: Patients are encouraged to continue performing regular self-breast examinations and to report any new or recurrent breast symptoms, including discharge, lumps, or skin changes, to their healthcare provider promptly.
  \item Activity Resumption: Gradual return to full physical activity is typically advised over 2-4 weeks, with specific guidance provided by the surgeon to avoid strenuous activities that could strain the incision.
  \item Warning Signs: Patients should be educated on warning signs that necessitate immediate contact with their surgeon, such as increasing redness, swelling, pain, fever, pus from the incision, or recurrence of nipple discharge.
\end{itemize}

Further imaging or specialist referral (e.g., to a medical oncologist or radiation oncologist) may be indicated if the pathology results reveal high-risk lesions or malignancy, or if new breast concerns arise during follow-up.

\section*{Outcomes \& Prognosis}
Nipple discharge is a common breast symptom, affecting approximately 5-10\% of women at some point in their lives. While most instances of nipple discharge are physiological (e.g., during pregnancy or lactation, or due to hormonal fluctuations), pathologic single-duct discharge, which is the indication for microdochectomy, is less frequent but clinically significant. This condition is most commonly observed in women between the ages of 40 and 60 years, though it can occur across a wide age range. 

Intraductal papilloma is by far the most frequent underlying cause, accounting for 50-80\% of cases of single-duct pathologic discharge. Other benign causes include duct ectasia and various forms of hyperplasia. Importantly, malignancy, including ductal carcinoma in situ (DCIS) or invasive carcinoma, is found in 5-15\% of cases of single-duct pathologic nipple discharge, with the incidence being higher for bloody or serosanguinous discharge. The procedure itself, microdochectomy, carries a very low mortality rate, primarily associated with general anesthesia risks. Morbidity is generally low and related to the specific risks outlined previously, such as hematoma, infection, or cosmetic changes, which are typically minor and manageable.

The outcomes following microdochectomy are generally excellent, particularly for symptom resolution and definitive diagnosis. The success rate for complete cessation of nipple discharge is very high, typically exceeding 90\%. The prognosis largely depends on the histopathological findings of the excised duct. For the majority of patients whose pathology reveals a benign condition, such as an intraductal papilloma without atypia, the prognosis is excellent, and the microdochectomy is considered curative for the discharge. 

Functional outcomes are typically good, with minimal long-term impact on breast function or cosmesis, and most patients report high satisfaction with the resolution of their symptoms. Recurrence of benign papillomas in the same or a different duct is possible but relatively uncommon, occurring in approximately 5-10\% of cases over time. For patients found to have atypical ductal hyperplasia (ADH) or ductal carcinoma in situ (DCIS) on microdochectomy, the prognosis is still generally favorable, but these findings necessitate further management, which may include re-excision to achieve clear margins, radiation therapy, or increased surveillance, depending on the specific diagnosis and institutional guidelines. 

Prognostic factors for long-term outcomes are primarily driven by the definitive pathology, the completeness of the excision, and the patient's overall risk profile for breast cancer. Landmark studies and guidelines from organizations like the American Society of Breast Surgeons consistently support microdochectomy as an effective and safe procedure for appropriate indications.

\section*{References}
\begin{enumerate}
  \item MedlinePlus. Microdochectomy. U.S. National Library of Medicine; 2025. Available from: \url{https://medlineplus.gov/ency/article/002931.htm}
  
  \item Schwartz's Principles of Surgery. 12th ed. Brunicardi FC, Andersen DK, Billiar TR, et al., editors. McGraw-Hill Education; 2023.
  
  \item Sabiston Textbook of Surgery: The Biological Basis of Modern Surgical Practice. 22nd ed. Townsend CM Jr, Beauchamp RD, Evers BM, Mattox KL, editors. Elsevier; 2021.
  
  \item American Society of Breast Surgeons. Consensus Guideline on the Management of Nipple Discharge. 2023. Available from: \url{https://www.breastsurgeons.org/statements/Nipple-Discharge-Guideline}
\end{enumerate}

\clearpage
